\section{Examples and the model contract}
\label{sec:examples}

All market examples use one frozen after-hours snapshot of SPY and NVDA listed
options from the 2026-08-03 US session.  Each underlier has eight expiries,
for \MacSnapNodeCount\ fitted nodes in total.  The common recipe is order
$N=16$ subject to \eqref{eq:orderguard}, the logistic chart, the haircut
target \eqref{eq:haircut}, ridge $\alpha=10^{-6}$, and the calendar rows
\eqref{eq:calrows}.  No node is tuned separately.  Across the sample the
median quote error is \MacSnapMedianRmsBp\ vol bp and the worst is
\MacSnapWorstRmsBp\ vol bp.

\subsection{SPY: a term structure and one slice}

\begin{figure}[htbp]
\centering
\paperfig{\textwidth}{fig_spy_gallery.pdf}
\caption{Eight SPY expiries fitted with one recipe.  Each panel shows raw
bid--ask quotes, mids, the fitted LQD slice, and rms error.  The short end is
steeper and more localized; the long end broadens smoothly.  Beyond the last
quote, each curve is determined by the fitted exponential-tail convention.}
\label{fig:spy_gallery}
\end{figure}

The gallery separates recurring market shape from the completion chosen by
the model.  Through the quoted region the slices follow the term structure of
level and skew; outside it the two endpoint speeds continue the curve.  The
boundary is most visible at the short end, where quoted spans are narrow in
log-moneyness: the sharp local skew is supported by quotes, while the remote
upward turn is the exponential-tail completion.  At longer maturities the
span widens and a larger fraction of the displayed smile is directly
informed by prices.  The median error across the eight slices is
\MacSpyGalleryMedianRmsBp\ vol bp; the worst is
\MacSpyGalleryWorstRmsBp\ vol bp.

\begin{figure}[htbp]
\centering
\paperfig{0.94\textwidth}{fig_spy_node.pdf}
\caption{The SPY December 2026 slice.  Left: quotes and fitted volatility.
Center: signed residuals in vol bp.  Right: the fitted density against an
ATM-matched normal.  The density view makes the negative skew visible as a
transfer of mass toward the left tail.}
\label{fig:spy_node}
\end{figure}

The December 2026 slice contains \MacSpyDecNQuotes\ quotes and fits them to
\MacSpyDecRmsBp\ vol bp rms.  Its exact ATM quantities from
\cref{prop:atm} are volatility \MacSpyDecAtmVolPct\%, skew
$\MacSpyDecSkew$, and curvature $\MacSpyDecCurv$.  The forward rank is
\MacSpyDecForwardRankPct\%, below the matched Black rank; hence
$\Delta<0$, in agreement with the negative skew.  The log-contract variance
corresponds to \MacSpyDecVarSwapPct\% annualized volatility.

The three panels answer different questions.  The first tests the nonlinear
map from law to implied volatility.  The second asks whether residuals retain
systematic strike structure.  The third removes the Black coordinate and
shows the fitted object itself: a left-heavy density.  The ATM-matched normal
has the same central price level but not the same digital or tail allocation,
and therefore cannot reproduce the skew.

The ledger also makes an individual option transparent.  For the listed
December \MacTicketStrike\ call, $k=\MacTicketK$ and the strike rank is
\MacTicketRankPct\%.  The upper-share and cash entries are
\MacTicketShare\ and \MacTicketCash, so
\[
c(\MacTicketK)=\MacTicketShare-\MacTicketCash=\MacTicketPrice.
\]
This is \MacTicketDollarPrice\ per share and corresponds to
\MacTicketIvPct\% implied volatility.

\subsection{NVDA and the order guard}

\begin{figure}[htbp]
\centering
\paperfig{0.92\textwidth}{fig_nvda_nodes.pdf}
\caption{Two NVDA slices.  Left: the shortest expiry, with
\MacNvdaOneDayNQuotes\ quotes and effective order
$N_{\rm eff}=\MacNvdaOneDayEffOrder$.  Right: the December 2027 expiry at
full order.  Wider NVDA spreads leave live haircut bands and permit visibly
larger residual variation than SPY.}
\label{fig:nvda_nodes}
\end{figure}

The shortest node fits at \MacNvdaOneDayRmsBp\ vol bp rms; the long node, with
\MacNvdaLongNQuotes\ quotes, at \MacNvdaLongRmsBp\ vol bp.  The latter has
heavier fitted tails than SPY.  In particular its left endpoint speed exceeds
one, which is admissible: only the right endpoint is constrained by the
finite-forward wall.  The consequence is a last finite inverse-moment order
\MacNvdaLongMomentLeft\ below one.

The shortest node is also the clearest example of the order guard.  A thin
strip cannot distinguish sixteen smooth modes, even though an optimizer can
assign values to them.  Reducing $N_{\rm eff}$ converts that lack of
information into visible residual uncertainty instead of invisible
oscillation between quotes.  The long node has enough observations to use the
full order and separate body curvature from tail continuation.

\subsection{What basis order asserts}

The synthetic benchmark is a martingale-normalized mixture of two lognormal
components.  It has a genuinely bimodal density but only gentle shoulder
structure in implied volatility.

\begin{figure}[htbp]
\centering
\paperfig{0.9\textwidth}{fig_doublehump.pdf}
\caption{The double-hump benchmark at $N=16$.  Left: the fitted smile matches
the synthetic quotes to \MacHumpFitRmsBp\ vol bp rms.  Right: the fitted
density recovers both modes, the trough, and the asymmetry of the target.}
\label{fig:doublehump}
\end{figure}

\begin{figure}[htbp]
\centering
\paperfig{0.9\textwidth}{fig_order_control.pdf}
\caption{The same target at orders 6 and 16.  Left: the implied-volatility
curves nearly coincide over the quoted region.  Right: the densities do not;
the lower-order fit displaces the modes and fills part of the trough.  Basis
order controls how much inter-quote structure the model is allowed to assert.}
\label{fig:order_control}
\end{figure}

At $N=16$ the fitted mode locations are \MacDhModeLeft\ and
\MacDhModeRight, against true locations \MacDhTrueModeLeft\ and
\MacDhTrueModeRight.  At $N=6$ both modes remain, but their shape is distorted.
The $L^1$ density error rises from \MacHumpLOneNSixteen\ to
\MacHumpLOneNSix, even though the two implied-volatility curves differ by at
most \MacHumpIvAgreeBp\ vol bp on the comparison grid.  Option quotes can be
nearly indifferent between laws that price density-sensitive structures
differently.

The lesson is not that high order is always preferable.  At $N=16$ the
synthetic target is known to contain fine structure, so extra modes recover a
feature known to be real.  In market data the truth is not supplied.  High
order then grants the model permission to assert finer inter-quote density
shape.  Digitals and narrow butterflies are sensitive to that permission;
broad vanilla prices may not be.  Basis order should be chosen for the
downstream instruments, not solely for the in-sample volatility residual.

\subsection{Position among existing models}

The construction joins several classical ideas.  Breeden--Litzenberger
identify option prices with a density \citep{BreedenLitzenberger1978}; Lee and
Benaim--Friz connect moments with wing slopes
\citep{Lee2004,BenaimFriz2009}; quantile modeling supplies the rank coordinate
\citep{Parzen1979,Gilchrist2000}; and Petersen--M\"uller use the same
log-quantile-density transform to place density-valued data in an
unconstrained function space \citep{PetersenMuller2016}.  The metalog family
is a related direct quantile expansion \citep{Keelin2016}.

SVI and SSVI work in volatility coordinates and characterize admissible
parameter regions \citep{Gatheral2004,GatheralJacquier2014}; arbitrage-free
smoothing works directly with constrained prices \citep{Fengler2009}.  LQD
instead makes the probability law primary.  Its contribution is the combined
architecture: structural one-slice validity, readable tail coordinates, and
one ledger for pricing, differentiation, density, variance, and calendar
order.

\subsection{The contract}
\label{sec:limits}

Two things remain: how to read a fitted slice, and what precisely is
claimed.

A fitted slice should be read in the same order as the theory.  First inspect quote residuals over the retained span.  Then read the
three ATM handles as local level, digital, and density information.  Next
inspect $g(u)$ and the density over the ranks occupied by quotes.  Only after
that should one read $\lambda_\pm$, the moment strip, and Lee slopes, since
these describe the extrapolated endpoints.  Finally compare the ledger with
neighboring expiries.  This sequence prevents a tail convention or a
calendar diagnostic from being mistaken for a directly observed quote.

\begin{table}[htbp]
\centering\small
\begin{tabular}{@{}>{\raggedright\arraybackslash}p{0.31\textwidth}
>{\raggedright\arraybackslash}p{0.31\textwidth}
>{\raggedright\arraybackslash}p{0.30\textwidth}@{}}
\toprule
\textbf{Structural} & \textbf{Numerically controlled} &
\textbf{Outside the claim}\\
\midrule
Positive density, unit mean, and butterfly-free calls for every admissible
slice & Interpolation monotonicity, traded-range density, and randomized
price audits & Atoms, mass at zero, density gaps, bounded support, and
Gaussian-or-faster tails\\[4pt]
Exact exponential tail scales, moment strip, and Lee limits & Tail
quadrature, exponent budgets, grid convergence, and finite-strike wing prices
& Tail parameters as observations, or asymptotic slopes as finite-strike
quotes\\[4pt]
Exact ATM handles, analytic Jacobian, and ledger characterization of calendar
order & Quote fit, basis-order control, multi-start checks, and calendar order
on common quoted support & Dynamics, jump decomposition, realized-variance
corrections, and calendar structure outside the controlled support\\
\bottomrule
\end{tabular}
\caption{The model contract: theorem, numerical control, and excluded scope.}
\label{tab:contract}
\end{table}

The central distinction is now visible in one place.  Probability supplies
validity.  Quotes supply information over a finite strip.  Basis, ridge,
tails, and calendar scope supply convention elsewhere.  The value of the LQD
representation is not that convention disappears, but that it can no longer
be mistaken for a theorem or for market information.
